\begin{tomtat}
	\begin{minipage}[t]{0.3\textwidth}
		\vspace{0pt}
		\begin{tomtat}
			\begin{center}
				\textbf{ĐỀ 18}
			\end{center}
			
		\end{tomtat}
	\end{minipage}
	\hfill
	\begin{minipage}[t]{0.5\textwidth}
		\vspace{0pt}
		\begin{tomtat}
			\begin{center}
				\includegraphics[scale=0.08]{images/18}
			\end{center}
		\end{tomtat}
	\end{minipage}
\end{tomtat}
\textbf{I. PHẦN TRẮC NGHIỆM}
\begin{ex} % Câu 1
	Giải phương trình: $(x+2)(x-1)=0$. Kết luận nào là đúng?
	\choice{$x=-2$}{$x=1$}{$x=-2$ và $x=1$}{$x=2$ và $x=1$}
	\loigiai{
		Ta có phương trình tích $(x+2)(x-1)=0$.
		Điều này xảy ra khi một trong hai thừa số bằng 0:
		\begin{itemize}
			\item $x+2=0 \Rightarrow x=-2$
			\item $x-1=0 \Rightarrow x=1$
		\end{itemize}
		Vậy nghiệm của phương trình là $x=-2$ và $x=1$.
		\\ \textbf{ĐÁP ÁN C}
	}
\end{ex}

\begin{ex} % Câu 2
	Bất phương trình nào sau đây là bất phương trình bậc nhất một ẩn?
	\choice{$-2x^2+5>0$}{$3x-y \le 0$}{$-4x-2<0$}{$5+0x \ge -7$}
	\loigiai{
		Bất phương trình bậc nhất một ẩn có dạng $ax+b<0$ (hoặc $>, \le, \ge$) với $a \ne 0$.
		\begin{itemize}
			\item A. $-2x^2+5>0$ là bất phương trình bậc hai một ẩn.
			\item B. $3x-y \le 0$ là bất phương trình bậc nhất hai ẩn.
			\item C. $-4x-2<0$. Đây là bất phương trình bậc nhất một ẩn với $a=-4, b=-2$.
			\item D. $5+0x \ge -7 \Leftrightarrow 5 \ge -7$ (luôn đúng). Đây không phải là bất phương trình bậc nhất một ẩn vì hệ số của $x$ là $0$.
		\end{itemize}
		Vậy bất phương trình $-4x-2<0$ là bất phương trình bậc nhất một ẩn.
		\\ \textbf{ĐÁP ÁN C}
	}
\end{ex}

\begin{ex} % Câu 3
	Kết quả thực hiện phép tính $\sqrt{2}-\sqrt{18}$ là?
	\choice{$-\sqrt{2}$}{$-2\sqrt{2}$}{$-3\sqrt{2}$}{$-4\sqrt{2}$}
	\loigiai{
		Ta có:
		\begin{align*}
			\sqrt{2}-\sqrt{18} &= \sqrt{2} - \sqrt{9 \cdot 2} \\
			&= \sqrt{2} - 3\sqrt{2} \\
			&= (1-3)\sqrt{2} \\
			&= -2\sqrt{2}
		\end{align*}
		\\ \textbf{ĐÁP ÁN B}
	}
\end{ex}

\begin{ex} % Câu 4
	Một tòa tháp có bóng trên mặt đất dài 15m, biết rằng góc tạo bởi tia nắng mặt trời với mặt đất là $55^\circ$. Chiều cao của tòa tháp (làm tròn đến hàng phần mười của mét) là:
	\choice{21,4}{8,6}{12,3}{10,5}
	\loigiai{
		Gọi $h$ là chiều cao của tòa tháp (cạnh đối) và $b = 15 \text{ m}$ là độ dài bóng trên mặt đất (cạnh kề).
		Góc tạo bởi tia nắng mặt trời và mặt đất là $\alpha = 55^\circ$.
		Ta có mối quan hệ: $\tan \alpha = \frac{\text{cạnh đối}}{\text{cạnh kề}} = \frac{h}{b}$.
		Do đó, $h = b \cdot \tan \alpha$.
		$$ h = 15 \cdot \tan(55^\circ) $$
		Sử dụng máy tính: $\tan(55^\circ) \approx 1,4281$.
		$$ h \approx 15 \cdot 1,4281 \approx 21,4215 \text{ m} $$
		Làm tròn đến hàng phần mười, ta được $h \approx 21,4 \text{ m}$.
		\\ \textbf{ĐÁP ÁN A}
	}
\end{ex}

\begin{ex} % Câu 5
	Cho đường tròn tâm O, chu vi 36$\pi$, cung có số đo $90^\circ$ có độ dài là:
	\choice{$18\pi$}{$4,5$}{$9\pi$}{$10\pi$}
	\loigiai{
		Chu vi của đường tròn là $C = 2\pi R = 36\pi$.
		Suy ra bán kính $R = \frac{36\pi}{2\pi} = 18$. (Đơn vị độ dài, ví dụ cm).
		Độ dài cung tròn có số đo $n^\circ$ được tính bằng công thức: $l = \frac{\pi R n}{180}$.
		Với $R=18$ và $n=90^\circ$:
		\begin{align*}
			l &= \frac{\pi \cdot 18 \cdot 90}{180} \\
			&= \frac{\pi \cdot 18 \cdot 1}{2} \\
			&= 9\pi
		\end{align*}
		\\ \textbf{ĐÁP ÁN C}
	}
\end{ex}

\begin{ex} % Câu 6
	Cho hàm số $y=\frac{1}{3}x^2$. Giá trị của $y$ tại $x=2\sqrt{3}$ là:
	\choice{$\frac{2\sqrt{3}}{3}$}{2}{$\frac{4\sqrt{3}}{3}$}{4}
	\loigiai{
		Thay $x=2\sqrt{3}$ vào hàm số $y=\frac{1}{3}x^2$:
		\begin{align*}
			y &= \frac{1}{3}(2\sqrt{3})^2 \\
			&= \frac{1}{3}(2^2 \cdot (\sqrt{3})^2) \\
			&= \frac{1}{3}(4 \cdot 3) \\
			&= \frac{1}{3}(12) \\
			&= 4
		\end{align*}
		Vậy giá trị của $y$ tại $x=2\sqrt{3}$ là 4.
		\\ \textbf{ĐÁP ÁN D}
	}
\end{ex}
\begin{ex} % Câu 7
	Thống kê thâm niên công tác (đơn vị: năm) của 33 nhân viên ở một công sở như sau:
	\begin{center}
		\begin{tabular}{|c|c|c|c|c|c|c|c|c|c|c|}
			\hline
			7 & 2 & 5 & 9 & 7 & 4 & 3 & 8 & 10 & 4 & 4 \\
			\hline
			2 & 4 & 4 & 5 & 6 & 7 & 7 & 5 & 4 & 1 & 8 \\
			\hline
			9 & 4 & 2 & 8 & 5 & 5 & 7 & 3 & 14 & 8 & 8 \\
			\hline
		\end{tabular}
	\end{center}
	Tần số của giá trị 2 năm là:
	\choice{2}{3}{4}{5}
	\loigiai{
		Đếm số lần xuất hiện của giá trị "2 năm" trong bảng dữ liệu:
		\begin{itemize}
			\item Dòng 1: 2 (1 lần)
			\item Dòng 2: 2 (2 lần)
			\item Dòng 3: 2 (3 lần)
		\end{itemize}
		Vậy, tần số của giá trị 2 năm là 3.
		\\ \textbf{ĐÁP ÁN B}
	}
\end{ex}

\begin{ex} % Câu 8
	Một hình trụ có bán kính đáy bằng 5cm, chiều cao bằng 8cm. Diện tích xung quanh hình trụ đó là:
	% Lưu ý: Đề bài ghi "hình lăng trụ" nhưng các dữ kiện (bán kính đáy, chiều cao) và các đáp án (có pi) gợi ý đây là hình trụ.
	\choice{$80\pi \text{ cm}^2$}{$200\pi \text{ cm}^2$}{$40\pi \text{ cm}^2$}{$100\pi \text{ cm}^2$}
	\loigiai{
		Diện tích xung quanh của hình trụ được tính bằng công thức: $S_{xq} = 2\pi Rh$.
		Trong đó:
		$R$ là bán kính đáy của hình trụ ($R=5 \text{ cm}$).
		$h$ là chiều cao của hình trụ ($h=8 \text{ cm}$).
		\begin{align*}
			S_{xq} &= 2\pi \cdot R \cdot h \\
			&= 2\pi \cdot 5 \cdot 8 \\
			&= 80\pi \text{ (cm}^2\text{)}
		\end{align*}
		\\ \textbf{ĐÁP ÁN A}
	}
\end{ex}
\newpage
\textbf{II. PHẦN TỰ LUẬN}
\begin{ex} % Câu 9
	Rút gọn các biểu thức $P=\left(\frac{1}{\sqrt{x}-1}-\frac{1}{\sqrt{x}}\right)\left(1-\frac{1}{\sqrt{x}}\right)$ với $x>0, x \ne 1$.
	\loigiai{
		Với $x>0, x \ne 1$, ta có:
		\begin{align*}
			P &= \left(\frac{\sqrt{x} - (\sqrt{x}-1)}{\sqrt{x}(\sqrt{x}-1)}\right) \cdot \left(\frac{\sqrt{x}-1}{\sqrt{x}}\right) \\
			&= \frac{\sqrt{x} - \sqrt{x} + 1}{\sqrt{x}(\sqrt{x}-1)} \cdot \frac{\sqrt{x}-1}{\sqrt{x}} \\
			&= \frac{1}{\sqrt{x}(\sqrt{x}-1)} \cdot \frac{\sqrt{x}-1}{\sqrt{x}}
		\end{align*}
		Vì $x \ne 1$ nên $\sqrt{x}-1 \ne 0$. Ta có thể rút gọn $\sqrt{x}-1$.
		\begin{align*}
			P &= \frac{1}{\sqrt{x}} \cdot \frac{1}{\sqrt{x}} \\
			&= \frac{1}{x}
		\end{align*}
		Vậy $P = \frac{1}{x}$.
	}
\end{ex}

\begin{ex} % Câu 10
	Gọi $x_1, x_2$ là nghiệm của phương trình $x^2-3x+1=0$. Không giải phương trình. Hãy tính giá trị của biểu thức $M=x_1^2+x_2^2-4x_1x_2$.
	\loigiai{
		Phương trình đã cho là $x^2-3x+1=0$.
		Đây là phương trình bậc hai có dạng $ax^2+bx+c=0$ với $a=1, b=-3, c=1$.
		Ta tính biệt thức $\Delta = b^2-4ac = (-3)^2 - 4(1)(1) = 9 - 4 = 5$.
		Vì $\Delta = 5 > 0$ nên phương trình có hai nghiệm phân biệt $x_1, x_2$.
		
		Theo hệ thức Vi-ét, ta có:
		$$ \begin{cases} 
			x_1+x_2 = \dfrac{-b}{a} = \dfrac{-(-3)}{1} = 3 \\ 
			x_1x_2 = \dfrac{c}{a} = \dfrac{1}{1} = 1 
		\end{cases} $$
		Ta cần tính giá trị của biểu thức $M=x_1^2+x_2^2-4x_1x_2$.
		Biến đổi $x_1^2+x_2^2 = (x_1+x_2)^2 - 2x_1x_2$.
		Vậy,
		\begin{align*}
			M &= (x_1+x_2)^2 - 2x_1x_2 - 4x_1x_2 \\
			&= (x_1+x_2)^2 - 6x_1x_2
		\end{align*}
		Thay giá trị $x_1+x_2 = 3$ và $x_1x_2 = 1$ vào biểu thức $M$:
		\begin{align*}
			M &= (3)^2 - 6(1) \\
			&= 9 - 6 \\
			&= 3
		\end{align*}
		Vậy giá trị của biểu thức $M$ là 3.
	}
\end{ex}

\begin{ex} % Câu 11
	Giải hệ phương trình $\begin{cases} 3x-2y=4 \\ 2x+y=5 \end{cases}$.
	\loigiai{
		Ta có hệ phương trình:
		$$ \begin{cases} 3x-2y=4 & (1) \\ 2x+y=5 & (2) \end{cases} $$
		Từ phương trình (2), suy ra $y = 5-2x$.
		Thay $y = 5-2x$ vào phương trình (1):
		\begin{align*}
			3x - 2(5-2x) &= 4 \\
			3x - 10 + 4x &= 4 \\
			7x &= 4 + 10 \\
			7x &= 14 \\
			x &= 2
		\end{align*}
		Thay $x=2$ vào $y = 5-2x$:
		\begin{align*}
			y &= 5 - 2(2) \\
			y &= 5 - 4 \\
			y &= 1
		\end{align*}
		Vậy nghiệm của hệ phương trình là $(x; y) = (2; 1)$.
	}
\end{ex}

\begin{ex} % Câu 12
	 Một đội xe theo kế hoạch phải chở 240 tấn hàng trong một thời gian dự định. Do điều kiện thời tiết không thuận lợi nên mỗi ngày đội chở ít hơn 6 tấn, do đó đã hoàn thành công việc chậm hơn 2 ngày so với dự định. Hỏi theo kế hoạch mỗi ngày đội xe phải chở bao nhiêu tấn hàng?
	\loigiai{
		Gọi số tấn hàng mỗi ngày đội xe dự định chở là $x$ (tấn/ngày).
		Điều kiện: $x > 6$.
		
		Thời gian dự định để chở hết 240 tấn hàng là: $\frac{240}{x}$ (ngày).
		
		Do điều kiện thời tiết không thuận lợi, mỗi ngày đội chở ít hơn 6 tấn, nên thực tế mỗi ngày đội chở được: $x-6$ (tấn/ngày).
		Thời gian thực tế để chở hết 240 tấn hàng là: $\frac{240}{x-6}$ (ngày).
		
		Vì đội hoàn thành công việc chậm hơn 2 ngày so với dự định, nên ta có phương trình:
		$$ \frac{240}{x-6} - \frac{240}{x} = 2 $$
		Quy đồng mẫu số:
		\begin{align*}
			\frac{240x - 240(x-6)}{x(x-6)} &= \frac{2x(x-6)}{x(x-6)} \\
			240x - (240x - 1440) &= 2x(x-6) \\
			240x - 240x + 1440 &= 2x^2 - 12x \\
			1440 &= 2x^2 - 12x \\
			2x^2 - 12x - 1440 &= 0
		\end{align*}
		Chia cả hai vế cho 2:
		$$ x^2 - 6x - 720 = 0 $$
		Tính biệt thức $\Delta' = (b')^2 - ac = (-3)^2 - 1(-720) = 9 + 720 = 729$.
		$\sqrt{\Delta'} = \sqrt{729} = 27$.
		Phương trình có hai nghiệm:
		$$ x_1 = \frac{-(-3) + 27}{1} = 3 + 27 = 30 $$
		$$ x_2 = \frac{-(-3) - 27}{1} = 3 - 27 = -24 $$
		Đối chiếu với điều kiện $x > 6$, ta nhận nghiệm $x=30$.
		
		Vậy theo kế hoạch mỗi ngày đội xe phải chở 30 tấn hàng.
	}
\end{ex}
\begin{ex} % Câu 12a (Từ ảnh mới)
	Một hộp chứa 1 viên bi xanh, 1 viên bi đỏ và 1 viên bi trắng. Các viên bi có cùng kích thước và khối lượng. An lần lượt lấy ra ngẫu nhiên từng viên bi từ trong hộp cho đến khi hết bi. Tính xác suất của biến cố sau: "Viên bi màu đỏ được lấy ra sau cùng".
	\loigiai{
		Ký hiệu các viên bi: Xanh (X), Đỏ (Đ), Trắng (T).
		An lần lượt lấy ra ngẫu nhiên từng viên bi cho đến khi hết bi. Đây là một hoán vị của 3 viên bi.
		
		Tập hợp tất cả các kết quả có thể xảy ra (không gian mẫu $\Omega$) là các thứ tự lấy bi: \\
		$\Omega = \{ $
		(X, Đ, T); (X, T, Đ); 
		(Đ, X, T); (Đ, T, X); 
		(T, X, Đ); (T, Đ, X) 
		$\}$.
		Số phần tử của không gian mẫu là $n(\Omega) = 6$.
		
		Gọi A là biến cố: "Viên bi màu đỏ được lấy ra sau cùng".
		Điều này có nghĩa là viên bi Đỏ (Đ) ở vị trí thứ ba.
		Tập hợp các kết quả thuận lợi cho biến cố A:\\
		$A = \{ $
		(X, T, Đ); (T, X, Đ)
		$\}$.
		Số phần tử của biến cố A là $n(A) = 2$.
		
		Xác suất của biến cố A là:
		$$ P(A) = \frac{n(A)}{n(\Omega)} = \frac{2}{6} = \frac{1}{3} $$
		Vậy xác suất để viên bi màu đỏ được lấy ra sau cùng là $\frac{1}{3}$.
	}
\end{ex}
\begin{ex} % Câu 13
	Cho tam giác $ABC$ vuông tại $A$. Lấy điểm $M$ bất kì trên đoạn $AC$, đường tròn đường kính $CM$ cắt hai đường thẳng $BM, BC$ lần lượt tại $D, N$.
	\begin{enumerate}
		\item[a)] Chứng minh bốn điểm $A, B, C, D$ cùng nằm trên một đường tròn.
		\item[b)] Chứng minh: $BM \cdot BD = BN \cdot BC$ và các đường thẳng $AB, MN, CD$ cùng đi qua một điểm.
	\end{enumerate}
	\loigiai{
	\begin{center}
	\begin{tikzpicture}[line join = round, line cap=round,>=stealth,font=\footnotesize,scale=1]
		% Def
		\def\r{3} % This was in your template, not used in this specific drawing directly
		
		% Coordinates
		\coordinate (A) at (0,0);
		\coordinate (B) at (0,5);
		\coordinate (C) at (4,0);
		\coordinate (T) at (0,-2);
		\coordinate (D) at (14/5, -3/5);  % (2.8, -0.6)
		\coordinate (N) at (140/41, 30/41); % approx (3.41, 0.73)
		\coordinate (M) at (5/2, 0);    % (2.5, 0)
		\coordinate (I) at (13/4, 0);   % (3.25, 0)
		
		% Draw
		% Draw circle (I, MC/2), radius is 0.75
		\draw (I) circle (0.75);
		
		% Draw B-T-C-B
		\draw (B) -- (T) -- (C) -- cycle;
		
		% Draw M-A-D-M-C-N-M-B
		\draw (M) -- (A) -- (D) -- (M) -- (C) -- (N) -- (M) -- (B);
		
		% Fill points
		\fill (A) circle (1.5pt);
		\fill (B) circle (1.5pt);
		\fill (C) circle (1.5pt);
		\fill (T) circle (1.5pt);
		\fill (D) circle (1.5pt);
		\fill (N) circle (1.5pt);
		\fill (M) circle (1.5pt);
		\fill (I) circle (1.5pt);
		
		% Node points
		\node[below left] at (A) {$A$};
		\node[above left] at (B) {$B$};
		\node[below right] at (C) {$C$};
		\node[below left] at (T) {$T$};
		\node[below right] at (D) {$D$};
		\node[above right] at (N) {$N$};
		\node[below] at (M) {$M$}; % M is on x-axis, below should be fine
		\node[below right] at (I) {$I$}; % I is on x-axis, below right
		
		% Ky hieu goc vuong
		\pic[draw, angle radius=3mm]{right angle = C--A--B}; % BAC
		\pic[draw, angle radius=3mm]{right angle = M--N--C}; % MNC
		\pic[draw, angle radius=3mm]{right angle = M--D--C}; % MDC
		
	\end{tikzpicture}	
	\end{center}
	}
\end{ex}

\begin{ex} % Câu 14
	\begin{enumerate}
		\item[a)] Một ô tô đang chuyển động trên đường thẳng $AC$ theo hướng từ $A$ đi về phía $C$ với vận tốc 10m/s, một người đứng tại $B$ cách mép đường một khoảng $BH=50\text{m}$. Khi khoảng cách giữa người và ô tô là $AB=200\text{m}$ thì người đó bắt đầu chạy ra đón ô tô (coi ô tô và người chuyển động thẳng đều). Tìm vận tốc tối thiểu và hướng chạy của người tạo với $AB$ góc bao nhiêu để đón được ô tô.
		\item[b)] Cho $a, b$ là các số thực dương thỏa mãn điều kiện: $a+b=4ab$. Chứng minh rằng:
		$$ \frac{a}{4b^2+1} + \frac{b}{4a^2+1} \ge \frac{1}{2}. $$
	\end{enumerate}
\end{ex}

\setcounterref{ex}{0}
